% Phasor-Face Transformers (PrismFormer) -- LaTeX edition.
% Companion to paper/prismformer.md; identical content, math typeset properly.
% Build: pdflatex -> biber -> pdflatex -> pdflatex  (biblatex; references live in prismformer.bib).
% On Overleaf: set the compiler to pdfLaTeX and it runs biber automatically.
\documentclass[11pt]{article}

\usepackage[T1]{fontenc}
\usepackage[utf8]{inputenc}
\usepackage{lmodern}
\usepackage[margin=1in]{geometry}
\usepackage{amsmath,amssymb}
\usepackage{booktabs}
\usepackage{array}
\usepackage{microtype}
\usepackage[hidelinks]{hyperref}
\usepackage{csquotes}
\usepackage[backend=biber,style=authoryear,sorting=nyt]{biblatex}
\addbibresource{prismformer.bib}

\DeclareMathOperator{\softmax}{softmax}
\DeclareMathOperator*{\argmax}{arg\,max}
\newcommand{\bind}{\mathbin{\circledast}}         % phasor bind (per-component complex multiply)
\newcommand{\face}[1]{\operatorname{face}(#1)}
\newcommand{\bank}{\operatorname{bank}}
\newcommand{\Emb}{\operatorname{Emb}}
\newcommand{\code}[1]{\texttt{#1}}

\setlength{\parindent}{0pt}
\setlength{\parskip}{0.6em}

\title{\textbf{Phasor-Face Transformers: Decodable, Generalising Arithmetic in a Small Transformer}\\[0.3em]
\large\normalfont\itshape PrismFormer: arithmetic as algebra, with bit-exact, mergeable gradients on plain CPUs.}
\author{\textbf{Dongyang Stephen Chen}\\ Evaluated Applications\\
\small\texttt{dongyang.stephen.chen@gmail.com}}
\date{Draft, 2026}

\begin{document}
\maketitle

\begin{center}
\textbf{DOI:} \href{https://doi.org/10.5281/zenodo.21384525}{10.5281/zenodo.21384525}\\[0.3em]
\small Cite as: Chen, D.\,S. (2026). \emph{Phasor-Face Transformers (PrismFormer): Decodable,
Generalising Arithmetic in a Small Transformer, with Bit-Exact, Mergeable Gradients.} Zenodo.
\href{https://doi.org/10.5281/zenodo.21384525}{https://doi.org/10.5281/zenodo.21384525}
\end{center}

\begin{center}
\small\itshape
The exact source, benchmarks, and runnable verification kit that produce every number below are
archived as a frozen snapshot alongside this document (repository tag \texttt{paper-v3}); every
quantitative result is produced by a script in that snapshot and reproducible with the command given
in its section.
\end{center}

\begin{abstract}
I present \textbf{PrismFormer}, a transformer-shaped sequence model in which every token, word or
number, is a \textbf{phasor face}: a bundle of unit complex numbers composed by one algebraic
operation, per-component complex multiplication (``binding''). Numbers are encoded so that binding
performs \emph{arithmetic}: addition on a linear-phase band, multiplication on a log-phase band, with
results read back by correlation, so a computed value is \emph{decodable} rather than merely reachable.
This numeric representation is a direct application of the frequency-domain (phasor) form of
\textbf{Holographic Reduced Representations} \parencite{plate1995hrr,plate2003hrr} and of \textbf{Random Fourier Features}
\parencite{rahimi2007random}; my contribution is to make it the trainable substrate of a transformer whose
dense maps are replaced by parameter-lean \textbf{relation-banks}. Because that substrate can be
evaluated read-only during backpropagation, gradients accumulate into a detached buffer: a minibatch
split into shards and summed is \textbf{bit-for-bit identical} to the serial result, which I verify to the
bit. That merge property is the foundation a \textbf{full-replication training swarm} would stand on,
every node holding the whole model on a plain CPU and staying coherent across mismatched hardware, but I
show it on one machine here, in process and over loopback sockets: the merged run converges
bit-identically to single-machine training (identical loss, accuracy and weights at every step), and skill
spreads without a coordinator, independently-initialised nodes share what they have learned through
continuous weight coupling, near-losslessly with the shared codec ($99\%$ versus $75\%$ cross-init
transfer). A geo-distributed run at scale is future work.

I report results honestly and reproducibly, against a parameter-matched dense transformer at each
setting. Arithmetic emerges from the codec with \textbf{no training} (binding two number faces decodes
to their exact sum and product; $100/100$ and $81/81$ single-digit pairs exact). Given the same
\emph{worked, column-by-column} problem, PrismFormer learns multi-digit addition, subtraction, and
multiply/divide by a digit and generalises to \textbf{unseen} operand pairs (add $54.9\%$, sub
$70.0\%$, mul $27.7\%$, div $33.2\%$ held-out) where a matched transformer, tuned hard (layer
normalization, warmup, tuned Adam), still stays near zero; and its working is legible one column at a time with no trained probe. On comparison
and relational tasks it generalises better than a transformer that \emph{fits the same training data}
(held-out $65.9\%$ vs $49.1\%$ on the fair tasks; both reach $100\%$ on copy). It is a competitive
character language model at matched size ($41.0\%$ vs $21.7\%$ next-character accuracy). I report a
two-sided ablation of my central design choice: freezing the numeric identity is \emph{not} what makes
any single operation generalise (with one task in isolation, an unfrozen identity does better), but
under realistic \textbf{shared multi-task load} freezing wins the held-out comparison, because the
identity is a residual every task shares, and drifting it to fit one operation corrupts the rest. So
freezing buys both legibility and shared-corpus generalisation, and I say exactly when each holds.
Out-of-range extrapolation and large-scale convergence of the full colony remain open; reverse
inference, $0\%$ for the closed-book trained model, is recovered by a scratchpad or an algebraic encoding.
\end{abstract}

\section{Introduction}

I should say up front what this is and where it came from, because the path is what shaped the design.
I am not a machine-learning researcher. I am a full-stack developer with a background in computer
networking, and this started as a stubborn question rather than a gap in a literature I did not yet
know: if you wanted a model to \emph{reason} instead of recall, what would you actually have to
represent, and how? The only instrument I had for that question was my own mind, so I used it as the
template, and tried as rigorously as I could to watch how I actually reason and ask what structure that
would demand.

I did not start from the literature. I started from a short list of near-philosophical axioms I will not
defend here (the paper does not rest on them) and, more usefully, one picture I borrowed and stripped
down: a \emph{Platonic space}, an idea I met through Michael Levin's writing
\parencite{levin_platonic_blog} and ancient long before either of us. I needed exactly one thing from it:
that a concept has a definite, deterministic \emph{address}, the number five, or a particular word, each
with one place it belongs, so that reasoning becomes moving between addresses under fixed rules rather
than looking up a table. In the codec that is true by construction, every token resolves to one address,
so every copy of the model agrees on where a concept lives, and that agreement is exactly what later lets
separate machines merge into one coherent model: they are all writing to the same coordinates.

Then I did the engineer's kind of research rather than the scientist's: I built the thing, tested it,
and mostly watched it fail, for about a year and across three separate systems. The first was pure
geometry with no neural network at all, concepts as positions and operations as fixed moves between
them. The second bolted a small learned \emph{navigator} onto that space. This paper is the third, where
the space is folded into the transformer itself.\footnote{The two earlier iterations are public:
\texttt{genesis-repl} (the pure-geometry substrate) and \texttt{genesis-nova} (substrate plus
navigator), both under \url{https://github.com/EvaluatedApplications}.} I wrote tests for every idea so I
could tell a real effect from wishful thinking, which is the scientific method even when what you are
after is a working artefact and not a theorem. It was Edison's thousand dead lightbulbs, not a
derivation. And when I looked back at what had survived all three iterations unchanged, it was always
one property: a \emph{homomorphism}. The setups that worked were the ones where combining two
representations mirrored combining the two things they stood for, the same structure going in and coming
out. So I kept only that and threw the rest away.

Keeping only that property is what forced the shape the rest of this paper is about: a \emph{codec}, a
fixed and deterministic way to write any concept, word or number, as a vector, chosen so that one
operation composes two of them and correlation reads the answer back. The encoding that fell out was
phasors: each concept a bundle of unit complex numbers, composed by complex multiplication. Write
numbers as phases and they \emph{add} on one band and \emph{multiply} on another, so arithmetic falls
out of the very same bind that composes words. That was the moment it felt real. I had not taught it to
add; the representation just did.

Only afterwards did I find that the mathematics I had backed into is decades old. Representing concepts
as unit-magnitude phasors and binding them by complex multiplication is the frequency-domain form of
Plate's Holographic Reduced Representations \parencite{plate1994thesis,plate1995hrr,plate2003hrr}, later called FHRR in the Vector
Symbolic Architectures literature \parencite{gayler2003vsa,kanerva2009hdc}; writing a scalar as phases of random
frequencies and reading it back by correlation is a Random Fourier Feature map \parencite{rahimi2007random}.
I cite that work in full in Section~\ref{sec:related} and claim none of it as mine. What I did was
reach it from the reasoning-space side instead of the kernel or VSA side, and then take it somewhere it
had not gone: make the phasor codec the trainable substrate of a transformer, and build around it the
training system a networking person would build.

That is the part where my actual background earns its keep, and it turns out to be the other half of
the paper. Because the substrate is only read, never mutated, while gradients pile up in a separate
buffer, a batch split across shards and summed lands on a result that is \emph{bit-for-bit identical}
to running it on one machine (Sections~\ref{sec:merge} and~\ref{sec:exactmerge}). Bit-exactness is a
networking property before it is a learning one: a slow node's contribution and a fast node's combine
into exactly the same bits, so a swarm of mismatched, unreliable CPUs \emph{could} train one coherent
model with no parameter server and no matched hardware, every node holding the whole model the way every
piece of a hologram still holds the whole image. I want to be clear that this is the architecture the
merge property enables, and that I demonstrate the property and simulate the coupling in a single process
here, not a live multi-node run. (That holographic picture of memory has a long history as a
\emph{metaphor} for the brain, in Pribram's holonomic theory \parencite{pribram1971languages,pribram1991brain}; I use it only as a
metaphor and lean on the mathematics of HRR, not on any neuroscience claim.)

For context on why numbers are worth this trouble: language models handle them unevenly. Standard
embeddings capture magnitude only partly, and that fidelity falls off with sub-word tokenization and
past modest magnitudes \parencite{wallace2019numeracy}, while numbers rarely get any special representational
treatment \parencite{thawani2021numbers}. Embedded as points along shared directions, two different magnitudes
look nearly collinear to a dot-product readout, so the model recalls a memorised answer instead of
computing a new one. The codec is a direct answer to that: a computed value is decodable, not merely
reachable.

So the paper is two halves of one idea, and I make three contributions:

\begin{enumerate}
\item \textbf{A phasor-face transformer substrate (Sections~\ref{sec:codec}--\ref{sec:algformer}).}
Every parameter and activation is a decodable face; one binding operation performs symbolic
association and arithmetic at once; numeric embeddings are seeded from the phasor codec so arithmetic
is decodable from the representation with no training; dense maps are replaced by relation-banks
costing $S\cdot d$ parameters instead of $d^2$.
\item \textbf{Exactly-mergeable gradients (Section~\ref{sec:merge}).} The substrate is read-only during
backprop, so a minibatch splits into independent shards whose summed gradients equal the serial result
\emph{bit for bit}. I verify this to the bit.
\item \textbf{The architecture this enables: a full-replication CPU swarm that shares skill
(Sections~\ref{sec:merge} and~\ref{sec:skillshare}).} Because every node would hold the whole model and
gradients merge exactly, such a swarm \emph{could} train one coherent model on mismatched CPUs, the
opposite design to sharded systems such as Petals \parencite{borzunov2023petals}, which split one model
across matched accelerators and lose it if a peer drops. And skill spreads without a coordinator: in a
single process that simulates the mesh coupling I show independently-initialised nodes sharing what they
have learned through continuous weight coupling, near-losslessly when they share the codec's coordinate
system ($99\%$ versus $75\%$). Every experiment here runs on one machine, in process or over loopback
sockets, so I present this as a proven merge-and-share property plus the architecture it enables, and flag
the geo-distributed run at scale as future
work (Section~\ref{sec:lim}).
\end{enumerate}

I try to be honest about scope, and I ablate my own design choices both ways: an ablation
(Section~\ref{sec:worked}) shows that \emph{freezing} the numeric identity does not help any single
operation in isolation (unfreezing does better there), yet it wins under shared multi-task load, so I
present it as legibility \emph{plus} a shared-corpus generalisation benefit, with the isolated-task
caveat stated. Large-scale convergence and out-of-range extrapolation are open problems I do not claim
to have solved.

\section{Related work and attribution}\label{sec:related}

I build directly on prior inventions and cite them plainly; the novelty here is in their combination
and in the systems consequence, not in the primitives. PrismFormer's core architecture, which I call an \textbf{AlgFormer} (an
\emph{algebraic} transformer), is in one line the transformer
block \parencite{vaswani2017attention} with three parts swapped in: \textbf{FHRR phasor binding} \parencite{plate1994thesis,plate2003hrr} as the compose operator, a \textbf{Random Fourier Feature} scalar code \parencite{rahimi2007random} as
the number embedding, and \textbf{banded circular operators} \parencite{cheng2015circulant,sindhwani2015structured}
in place of every dense matrix, trained by data-parallel SGD \parencite{mcmahan2017fedavg} made bit-exact. The
rest of this section credits each part; Section~\ref{sec:method} shows how they compose into a whole
model.

\textbf{Binding and holographic representations.} Tensor-product binding \parencite{smolensky1990tensor} and
Holographic Reduced Representations \parencite{plate1995hrr} introduced fixed-width distributed representations of
structure with a circular-convolution bind and a correlation unbind. The frequency-domain form (Plate's
``unitary vectors'' \parencite{plate1994thesis,plate2003hrr}, later named FHRR in the VSA literature) uses unit-magnitude
\textbf{phasors} bound by element-wise complex multiplication, the exact bind I use. Vector Symbolic
Architectures \parencite{gayler2003vsa} and Hyperdimensional Computing \parencite{kanerva2009hdc}, surveyed by \textcite{kleyko2022vsa1,kleyko2023vsa2}, generalise this family. My phasor codec is an instance of this frequency-domain HRR; I do
not claim to have invented phasor binding.

\textbf{Encoding scalars as phases.} Writing a value $v$ as $(\cos v\theta_k,\ \sin v\theta_k)$ over
frequencies $\theta_k$ and reading it back by correlation is a Random Fourier Feature map \parencite{rahimi2007random}: the correlation approximates a kernel peaked at the true value, which is why a computed
result is \emph{decodable} rather than merely reachable. Prior work characterises how numeracy is only
partially captured by standard embeddings and where it breaks down \parencite{wallace2019numeracy}, and surveys
the broader need for better numeric representation \parencite{thawani2021numbers}.

\textbf{Numbers in transformers, and where I sit in that cluster.} Since I started this I have found a
line of recent work aimed squarely at getting numbers into transformers, and I want to place myself in
it honestly rather than let a reader think I missed it. Several of these overlap with things I claim, so
I will concede the overlaps first and only then draw the line I think is real. \textbf{xVal}
\parencite{golkar2023xval} tokenises a number as a single embedding scaled by its value, and \textbf{Abacus}
\parencite{mcleish2024abacus} adds position embeddings over digits so a transformer can carry arithmetic
across columns; both improve how a standard transformer handles numbers without changing what the number
\emph{is} representationally. Closer to me, \textbf{FoNE} \parencite{zhou2025fone} embeds a number directly with
Fourier features and \emph{decodes the digits back out by correlation}, so ``a computed number is
decodable, not just reachable'' is \textbf{not} unique to my codec, and I concede that plainly. Closer
still, \textbf{Neural Isomorphic Fields} (NIF) \parencite{sadeghi2026nif} explicitly builds an embedding in
which $\phi(a+b)=\phi(a)\,\dot{+}\,\phi(b)$ and $\phi(a\times b)=\phi(a)\,\dot{\times}\,\phi(b)$, so
``algebra in the embedding'' is not first either, and I concede that too.

Having conceded both, here is the distinction I do stand on. NIF \emph{learns} two separate operators
$\dot{+}$ and $\dot{\times}$ (plus an ordering network) and fits them with an isomorphism loss; its
multiplication is reported around $53$--$73\%$ and reads to me as extrapolation rather than an exact
identity. In PrismFormer there is \textbf{one} fixed operation, per-component complex multiply, and add
and multiply are the \emph{same} bind acting on two differently-encoded bands: a linear-phase band where
that bind adds values and a log-phase band where the identical bind multiplies them. Multiplication is
therefore \textbf{native}, the same additive move on the log band, not the weak spot and not a learned
second operator. And the same bind that composes numbers composes \emph{words}: FoNE and NIF are
number-only modules, whereas my codec puts words and numbers in one space under one algebra. Holographic
binding inside a transformer is itself prior art, \textbf{Hrrformer} \parencite{alam2023hrrformer} recasts
self-attention with HRR bind/unbind, but it spends the binding on attention \emph{efficiency} ($O(n\log
n)$ over long sequences), not on numeric semantics, so I cite it as the precedent for holographic
binding-in-a-transformer aimed at a different target. Two ingredients have no counterpart in any of these:
the parameter-lean \textbf{relation-banks} that replace every dense map, and the
\textbf{exactly-mergeable gradients} (a detached buffer summed in fixed order gives bit-identical sharded
sums), which is an orthogonal distributed-training property, not a numeric one. So the novelty I actually
claim is the \emph{combination}, one fixed complex-multiply bind, over linear and log bands, decodable by
correlation, unified across words and numbers, plus relation-banks and exact-merge, and \textbf{not}
``decodable numbers'' or ``algebra in the embedding'' taken alone, since those already exist.

\textbf{Parameter-efficient linear maps.} Replacing dense $W\!\cdot\! x$ with structured operators is
well studied: circulant projections \parencite{cheng2015circulant} and structured transforms \parencite{sindhwani2015structured}. My relation-bank $y[i] = \sum_{k<S} \bank[k][i]\cdot x[(i+k)\bmod d]$ is a sum of $S$
diagonally-scaled cyclic shifts (a ``banded'' family over the cyclic group $\mathbb{Z}/d\mathbb{Z}$);
at $S=d$ it reparametrises a full matrix, so $S$ is a lean-to-full-rank control.

\textbf{Architecture and optimisation.} The block structure (attention, then per-token feed-forward,
with residuals) follows the transformer \parencite{vaswani2017attention}; I train with Adam \parencite{kingma2015adam}.

\textbf{Distributed and decentralised training.} Data-parallel synchronous SGD sums gradients across
workers; my exact-merge is that identity made bit-exact on a CPU substrate. Weight-slice averaging
between peers is a form of elastic/federated averaging \parencite{mcmahan2017fedavg}. In contrast to
model-parallel collaborative inference \parencite{borzunov2023petals}, every PrismFormer node replicates the
whole model, so no peer is load-bearing.

\section{Method}\label{sec:method}

\subsection{The phasor codec}\label{sec:codec}

A \textbf{face} is a vector of $C = 128$ unit complex numbers, stored as $2C = 256$ interleaved reals
$(\cos\varphi_0,\ \sin\varphi_0,\ \cos\varphi_1,\ \sin\varphi_1,\ \dots)$. Write it as
$f = (e^{i\varphi_0}, \dots, e^{i\varphi_{C-1}})$; the meaning lives in the phases $\varphi$. Three
operations act on faces:

\begin{itemize}
\item \textbf{bind} $\bind$, per-component complex multiplication $(a \bind b)_c = a_c b_c$, so phases
\textbf{add} per component, $\varphi^a_c + \varphi^b_c$. In the interleaved reals:
\[
h[2c] = a[2c]\,b[2c] - a[2c{+}1]\,b[2c{+}1], \qquad
h[2c{+}1] = a[2c]\,b[2c{+}1] + a[2c{+}1]\,b[2c].
\]
\item \textbf{unbind}, bind with the conjugate, so phases \textbf{subtract}.
\item \textbf{bundle}, component-wise real sum, superposing faces into a set.
\end{itemize}

Readout is \textbf{correlation}, the real inner product
$\langle a, b\rangle = \sum_c \big(a[2c]\,b[2c] + a[2c{+}1]\,b[2c{+}1]\big)$, optionally restricted to a
band of components. The $C$ components split into three bands:

\begin{center}
\begin{tabular}{@{}llccc@{}}
\toprule
band & comps & phase written for value $v$ & what $\bind$ does & frozen \\
\midrule
linear  & 16 & $\varphi_k = v\,\theta_k$              & \textbf{adds} values       & yes (identity) \\
log     & 16 & $\varphi_k = \ln v\,\theta_k$ ($v>0$)  & \textbf{multiplies} values & yes (identity) \\
orbital & 96 & learned tail (symbol seed / trained)  & ---                        & no (learned)   \\
\bottomrule
\end{tabular}
\end{center}

The frequencies $\theta_k$ are fixed and deterministic (drawn from a hash). The bind identities are then
just exponent arithmetic. For two number faces on the \textbf{linear} band,
$e^{i v \theta_k}\, e^{i w \theta_k} = e^{i (v+w)\theta_k}$, so the bound face \textbf{is} the number
face of $v+w$; on the \textbf{log} band,
$e^{i \ln v\,\theta_k}\, e^{i \ln w\,\theta_k} = e^{i \ln(vw)\,\theta_k}$, the number face of $v\cdot w$.
One bind therefore computes a sum on one band and a product on the other at the same time. The value is
read back by correlation against the codebook:
$\mathrm{DecodeSum}(h) = \argmax_c \langle \face{c}, h\rangle$ over the linear band, and
$\mathrm{DecodeProduct}$ the same over the log band. As a Random Fourier Feature map \parencite{rahimi2007random}, that correlation is a kernel peaked at the true value, so a \emph{computed} result is nameable,
not merely reachable: there is no table and no arithmetic unit, only a phase pattern named by
correlation.

A non-numeric token is handled by the same machinery with a different on-ramp: it hashes (FNV-1a) to a
deterministic phasor \textbf{signature} on the identity band and a \textbf{seed} on the orbital tail.
Numbers and symbols thus share one space and one algebra; the only difference is that a number's
\emph{value} becomes phase while a symbol's \emph{spelling} does. The identity band (value or signature)
is a concept's ground truth and is \textbf{frozen}; the orbital tail is what training moves. This is
exactly the frequency-domain HRR of \textcite{plate1994thesis,plate2003hrr} with an RFF scalar encoding \parencite{rahimi2007random}; the contribution (Section~\ref{sec:algformer}) is to make it a trainable transformer substrate.

\emph{Verifiable now, no training:} $\bind\!\big(\face{6}, \face{9}\big)$ decodes to $15$ (linear band)
and $54$ (log band). Over all single-digit pairs the codec is exact: $\mathbf{100/100}$ sums (0--9) and
$\mathbf{81/81}$ products (1--9). Run \code{dotnet run --project verify/Verify -c Release}.

\subsection{The relation-bank transformer (AlgFormer)}\label{sec:algformer}

AlgFormer keeps the transformer block shape \parencite{vaswani2017attention} and changes two things: every
activation is a face, and every dense linear map $W\!\cdot\! x$ (a $d\times d$ matrix, $d^2$ parameters)
becomes a \textbf{relation-bank} of $S$ faces ($S\cdot d$ parameters):
\begin{equation}
y[i] \;=\; \sum_{k=0}^{S-1} \bank[k][i]\cdot x[(i+k)\bmod d].
\label{eq:bank}
\end{equation}
Read it algebraically: $x[(i+k)\bmod d]$ is $x$ cyclically rotated by $k$, $\bank[k]$ scales it
component-wise, and the sum over $k$ bundles $S$ such shifted-and-scaled copies. So a bank is a sum of
$S$ diagonally-scaled cyclic shifts, a \emph{banded circular operator} over $\mathbb{Z}/d\mathbb{Z}$. At
$S=d$ the shifts span the group and the bank reparametrises a full $d\times d$ matrix; at $S\ll d$ it is
lean. $S$ is the dial between the two, and I run it well below $d$ (single digits to a few dozen shifts
against $d = 256$).

A block takes the $T\times d$ stack $X$ of input faces (row $X[t]$ = the face at position $t$) and
computes:
\begin{enumerate}
\item \textbf{Embed.} $X[t] = \Emb[\mathrm{tok}_t] + \mathrm{Pos}[t]$. $\Emb$ is seeded from the codec
(Section~\ref{sec:codec}) with its identity band frozen; $\mathrm{Pos}$ is a learned position face.
\item \textbf{Project.} $Q = \bank_Q(X)$, $K = \bank_K(X)$, $V = \bank_V(X)$: three separate
relation-banks.
\item \textbf{Causal attention.}
$\displaystyle a[t][j] = \softmax_{j\le t}\!\left(\frac{\langle Q[t], K[j]\rangle}{\sqrt{d}}\right)$,
then $\displaystyle \mathrm{ctx}[t] = \sum_{j\le t} a[t][j]\,V[j]$. Position $t$ attends only to
$j\le t$, so the same forward pass serves training and KV-cached generation.
\item \textbf{Attention residual.} $O = \bank_O(\mathrm{ctx})$, $M = X + O$.
\item \textbf{Gated feed-forward.} $A = \bank_A(M)$, $G = \bank_G(M)$, then $Z = A \odot \sigma(G)$ (a
GLU gate; a complex-bind variant $Z = A \bind G$ is also gradient-checked), and $F = \bank_F(Z)$.
\item \textbf{FFN residual.} $Y = M + F$, passed to the next block. Depth is $L$ blocks.
\end{enumerate}

There is \textbf{no layer normalization} anywhere: the residual stream and the bounded gate keep
activations in range, and a linear learning-rate decay keeps the multiplicative feed-forward stable. The
readout is \textbf{tied} to the embedding table,
$\mathrm{logit}_w = C_w + \langle \Emb[w], Y_{\text{last}}\rangle$, so the model names its answer by
correlating the final face against every token's face, the \emph{same} correlation that decodes the
codec (Section~\ref{sec:codec}). Training minimises cross-entropy on that softmax. The whole backward
pass is hand-derived and finite-difference gradient-checked at production dimensions; the inner cell is
SIMD-vectorised over the output index (\code{System.Numerics.Vector<double>}), bit-identical to the
scalar form.

Each face separates a fixed \textbf{identity} region (the numeric value / symbol signature, a fixed-size
prefix of each face's reals) from a learnable \textbf{orbital} tail. My implementation \emph{freezes}
the identity region during training. I treat this as a design choice and test it both ways
(Section~\ref{sec:worked}): freezing does \textbf{not} improve generalisation on a single operation in
isolation (an unfrozen, still codec-seeded identity does better there), but it \textbf{wins} under
shared multi-task load, where the identity is a residual every task shares and drifting it to fit one
operation hurts the others. On top of that it buys legibility: the identity stays on the codec geometry,
so activations remain decodable with no probe (Section~\ref{sec:inside}). The load-bearing ingredient is
the codec-seeded initialisation; freezing then adds shared-corpus robustness and readability.

\subsection{Exactly-mergeable gradients and the swarm}\label{sec:merge}

The loss over a batch is a sum of per-example losses, $L = \sum_i \ell_i$, so the gradient is the sum of
per-example gradients, $\nabla L = \sum_i \nabla \ell_i$. That much is ordinary. Two implementation
facts make the sum reproduce \textbf{bit for bit} regardless of how it is scheduled:
\begin{enumerate}
\item \textbf{The substrate is read-only during backprop.} The forward pass caches activations; the
backward pass writes only into a \emph{detached} gradient buffer and thread-local scratch, mutating no
parameter or shared state. So $\nabla \ell_i$ is a pure, deterministic function of example $i$ alone;
shards share nothing.
\item \textbf{The reduction is fixed-order.} Gradient buffers are summed in a fixed parameter-index
order, and the algebraic cell itself accumulates in ascending-shift order. Floating-point addition is
not associative, so the \emph{order} is what fixes the bits; fixing it makes the result independent of
thread timing.
\end{enumerate}

Put together: split a batch into shards, accumulate each shard's $\nabla \ell_i$ into its own buffer,
sum the buffers in the fixed order, take one optimiser step. Running the shards \textbf{in parallel
gives a gradient identical to the bit to running them serially}, verified by an FNV checksum over every
gradient value (Section~\ref{sec:exactmerge}). There is no computation graph to synchronise and no
parameter server. (Freezing is one clearing of the identity-band gradient before the step, so exact
numbers never drift.)

This is what the swarm is built on. Every node holds the whole model (full replication), trains on its
own data, and contributes gradient buffers (or averaged weight slices \parencite{mcmahan2017fedavg}) that merge
without a coordinator. Because a slow CPU's buffer and a fast CPU's buffer combine into the same bits, a
colony of mismatched, unreliable machines stays one coherent model, and a new node joins simply by
receiving the model and, from then on, replicating and spreading it. That is the sense in which the
swarm is \emph{self-replicating}. Capacity grows in place: extending the shift count or context
zero-pads the new parameters, which contribute nothing to the algebraic sum until trained, so a
checkpoint upgrades without a retrain. What I show here is that the mechanism runs and merges exactly
(Section~\ref{sec:exactmerge}); convergence of a large live colony is left open (Section~\ref{sec:lim}).

\section{Experiments}

\subsection{Setup}

I compare PrismFormer against a \textbf{parameter-matched} dense transformer (a real encoder-style
transformer with hand-derived, gradient-checked backprop) whose shape is auto-searched to fit the same
budget at a trainable depth. Both models share vocab, context, training order, LR schedule, and
parameter budget at each setting; both pass a startup gradient check. Every number below is a mean over
several random seeds with its standard deviation, produced by the listed command.

\paragraph{Confirmatory versus exploratory, and what I do and do not claim statistically.} I did not
pre-register anything, so I label rather than pretend. I treat as \textbf{confirmatory} the matched-budget
head-to-heads that test the paper's headline claims directly: the worked multi-digit arithmetic gap
(Section~\ref{sec:worked}), the relational-generalisation aggregate (Section~\ref{sec:relational}), the
character language-model comparison (\S\,4.5), the codec-seed confound control (Section~\ref{sec:codecbaseline}),
and the exact-merge check (Section~\ref{sec:exactmerge}), which is deterministic and not statistical at
all. I treat as \textbf{exploratory} the diagnostics that probe \emph{why}, where I read the direction of
an effect rather than a calibrated margin: the freeze-versus-unfreeze ablation in both regimes
(Section~\ref{sec:worked}), the skill-sharing coupling arms (Section~\ref{sec:skillshare}), and the
reverse-inference recovery arms (Section~\ref{sec:revinfer}). On significance I am honest that I report
means and standard deviations over 4--8 seeds, not confidence intervals or hypothesis tests, so a reader
should take the SDs as spread and not as inference. Where I lean on a comparison I check the coarse thing
I can check without raw per-seed data, whether the gap exceeds the two means' combined spread: it does,
by a wide margin, for the confirmatory arithmetic gap (e.g.\ add $54.9\%\pm8.5\%$ vs $0.4\%\pm0.6\%$) and
the language model ($41.0\%\pm1.9\%$ vs $21.7\%\pm3.6\%$), and it does \emph{not} for the AlgFormer-vs-random-transformer
row of the codec-seed control ($45.8\%\pm15.9\%$ vs $31.3\%\pm21.3\%$), which is exactly why I read that
table only for the per-seed-consistent ``seeding hurts the transformer'' effect and not for an
architectural margin. Where the spread swallows the gap I say so at the table rather than here.

\begin{center}
\begin{tabular}{@{}lrrll@{}}
\toprule
experiment (command) & PrismFormer & transformer & seeds $\times$ budget & data / held-out \\
\midrule
\S\,4.2 comparison/relational (\code{bench}) & 295{,}936 & 297{,}792 & $5\times150$ ep & 1{,}112 train / 309 held \\
\S\,4.3 worked arithmetic (\code{--columnar}) & 353{,}296 & 352{,}784 & $4\times15$k steps & operands 0--99, $\sim$20\% held \\
\S\,4.4 mechanistic (\code{--inspect}) & 177{,}421 & 169{,}213 & $8\times6$k steps & single-digit, $\sim$20\% held \\
\S\,4.5 language (\code{--lm}) & 249{,}920 & 259{,}776 & $5\times30$ ep & 1{,}670-char corpus \\
\S\,4.6 exact-merge (\code{verify/Verify}) & \code{Mini} & --- & deterministic & --- \\
\S\,4.7 codec-seed control (\code{--codec-baseline}) & 241{,}696 & 804{,}896 & $5\times800$ ep & add/sub 0--12, $\sim$20\% held \\
\S\,4.8 skill-sharing (\code{--average}) & \code{Mini} & --- & elastic coupling & add 0--12, cross-init \\
\S\,4.9 reverse inference (\code{--revinfer}) & \code{Mini} / algebra & --- & 3 arms & 24 symmetric pairs \\
\bottomrule
\end{tabular}
\end{center}

\subsection{Comparison and relational generalisation}\label{sec:relational}

On tasks where the answer is \textbf{not} a trivial algebraic function of the input faces, so the codec
gives no shortcut, I test whether each model learns the relation and generalises to a held-out split
(20\% of instances by a fixed hash, never trained). I report per-task \textbf{train and held-out}
accuracy, so the transformer's fit is visible: it is a competent baseline, not a strawman.

\begin{center}
\begin{tabular}{@{}lcc@{}}
\toprule
task & transformer (train / held) & PrismFormer (train / held) \\
\midrule
copy (first of three) & $100\%$ / $\mathbf{100\% \pm 0\%}$ & $100\%$ / $\mathbf{100\% \pm 0\%}$ \\
greater-than & $77\%$ / $65\% \pm 9\%$  & $100\%$ / $\mathbf{97\% \pm 3\%}$ \\
max          & $84\%$ / $66\% \pm 23\%$ & $100\%$ / $\mathbf{98\% \pm 3\%}$ \\
min          & $75\%$ / $61\% \pm 28\%$ & $100\%$ / $\mathbf{98\% \pm 2\%}$ \\
parity       & $77\%$ / $53\% \pm 14\%$ & $100\%$ / $\mathbf{68\% \pm 11\%}$ \\
\bottomrule
\end{tabular}
\end{center}

The transformer \textbf{fits its training data} ($85.8\%$ mean train accuracy on these tasks) and
\textbf{ties} PrismFormer at $100\%$ on copy, so the comparison is fair. Across the \textbf{seven} fair tasks (the five above plus \emph{antonym} and \emph{capital})
PrismFormer generalises markedly better: $\mathbf{65.9\%}$ vs $\mathbf{49.1\%}$ held-out, averaged over
tasks and seeds. Tuning the baseline hard (pre-norm layer normalization, a warmup schedule and a tuned
optimiser; Section~\ref{sec:lim}) lifts it only to $50.8\%$, so the gap is not a weak-baseline artefact.
Two honest caveats. The two fair tasks not shown, antonym and capital, are tested only
in the held-out \emph{reverse} direction (train ``opposite of A = B'', hold out ``opposite of B = A'';
train ``capital of C = city'', hold out the inverse), and there \textbf{both models score
$\approx 0\%$}: neither infers the symmetric relation at this scale. Those two zeros are \emph{included}
in the $65.9$/$49.1$ mean, which is why it sits below the individual relation accuracies above. I flag
this \emph{reverse-inference} as a limitation (Section~\ref{sec:lim}) rather than dropping it from the
average, though the codec's own machinery recovers it with a scratchpad or an algebraic encoding
(Section~\ref{sec:revinfer}). And single-token arithmetic (add/sub/mul/div where the whole answer is one
vocab token) is \emph{representation-favoured}: the codec makes the answer face equal to
$\bind(\text{inputs})$, so PrismFormer decodes it by construction (held-out 27--64\% vs the
transformer's 3--21\%). That is a property of the encoding, not evidence a transformer cannot add; the
fair arithmetic test is worked, multi-digit, and appears next. Run \code{dotnet run --project bench -c
Release}.

\subsection{Worked multi-digit arithmetic: the algorithm, not the table}\label{sec:worked}

A sequence model cannot emit an exact multi-digit result in one token, so the fair test is to hand both
models the same \emph{worked} problem and see who learns the algorithm. I write each answer digit by
digit, least significant first, so carries flow causally, and the model must generate those digits
itself and carry across columns; no outside loop does the work. Each of the four operations trains its
own model at full budget (all four, plus their transformer baselines, in parallel), on operands over
0--99, with about 20\% of pairs held out by a fixed hash. I use the atomic single-pass columnar form of
each operation: add/sub are two-operand (carry / borrow); multiply is a number times a single digit and
divide a number by a single digit, so every answer digit is one local column step: a digit-multiply on
the log-phase band, an add or subtract on the linear band. A seen-distribution control tells a memorised
table (held-out far below seen) from a learned algorithm (held-out about equal to seen).

\begin{center}
\begin{tabular}{@{}lccc@{}}
\toprule
op & PrismFormer held-out [seen] & transformer (tuned) held-out [seen] & face-decode \\
\midrule
add & $\mathbf{54.9\% \pm 8.5\%}$ [58\%] & $0.6\% \pm 0.4\%$ [1\%] & 69\% \\
sub & $\mathbf{70.0\% \pm 8.9\%}$ [74\%] & $2.3\% \pm 1.1\%$ [2\%] & 67\% \\
mul & $\mathbf{27.7\% \pm 3.3\%}$ [43\%] & $11.5\% \pm 3.9\%$ [9\%] & 61\% \\
div & $\mathbf{33.2\% \pm 12.9\%}$ [43\%] & $2.8\% \pm 1.5\%$ [4\%] & 50\% \\
\bottomrule
\end{tabular}
\end{center}

Here \emph{face-decode} is the answer read directly off the frozen identity bands of the final computed
face, by correlation and with no trained probe (the legibility measure of Section~\ref{sec:inside});
the accuracy column, by contrast, is the token the trained readout actually \emph{emits}. That
face-decode sits \emph{above} accuracy (e.g.\ $69\%$ vs $54.9\%$ on add, $61\%$ vs $28\%$ on mul) is not a
glitch, and it is the point: the algebra has already placed the correct value in the representation more
often than the argmax readout emits the matching token, so the gap between the two columns is a lossy
final projection sitting on top of a more-correct computation, not a failure of the arithmetic itself.
PrismFormer computes all four; the transformer essentially none, and this is a \emph{tuned} baseline
(pre-norm LayerNorm with a gradient-checked backward, a warmup$\to$cosine schedule and tuned Adam,
parameter-matched): tuning does not rescue the task. It stays near zero even on pairs it \emph{was}
trained on for add, subtract and divide, and only partially learns multiply-by-a-digit ($11.5\%$ held-out,
about its $9\%$ seen), while PrismFormer answers unseen problems far better than the transformer answers
ones it has already seen. For add/sub, held-out matches seen (55 against 58, 70 against 74): it learned the column
algorithm, not a lookup. Multiply and divide sit below their seen numbers (28 against 43, 33 against
43); with a single-digit second operand the pair space is smaller, so part of the score is
memorisation, but 28--33\% on genuinely unseen pairs, well above the tuned transformer's $2.8$--$11.5\%$,
is still the column step generalising. Full multi-digit $\times$ multi-digit multiplication is left to future work:
producing the product digits directly is a digit \emph{convolution}, not a single-pass column, so it
needs an explicit partial-product scratchpad (single-digit multiplies on the log band, summed on the
linear band), the same two primitives, composed.

\paragraph{Ablation: does freezing the numeric identity help?} I hypothesised that freezing the
identity region (keeping numbers on the exact codec geometry) would aid generalisation. The answer
depends on the training regime, and the two regimes tell a consistent story. Trained on \textbf{one
operation in isolation} (this section's setup), the identity made \emph{learnable}
(\code{frozenPrefix=0}, still codec-seeded) generalises \textbf{better} on every op:

\begin{center}
\begin{tabular}{@{}lcc@{}}
\toprule
op & frozen identity & learnable identity \\
\midrule
add & $54.9\% \pm 8.5\%$  & $\mathbf{70.9\% \pm 13.6\%}$ \\
sub & $70.0\% \pm 8.9\%$  & $\mathbf{73.4\% \pm 8.0\%}$ \\
mul & $27.7\% \pm 3.3\%$  & $\mathbf{36.7\% \pm 8.0\%}$ \\
div & $33.2\% \pm 12.9\%$ & $\mathbf{40.8\% \pm 3.2\%}$ \\
\bottomrule
\end{tabular}
\end{center}

This is what you would expect: with a single task and nothing else to protect, an adaptable
pass-through is simply extra capacity the model folds that operation's transform into. But the identity
is a \emph{shared} residual: the same component carries every token's own value through every task.
Re-running the ablation under \textbf{shared multi-task load} (Section~\ref{sec:relational}'s corpus:
seven relational tasks plus all four arithmetic ops, one model, 5 seeds) reverses the result. Held-out
accuracy, mean over seeds:

\begin{center}
\begin{tabular}{@{}lcc@{}}
\toprule
held-out group & frozen identity & learnable identity \\
\midrule
relational tasks & $\mathbf{65.9\%}$ & $65.2\%$ \\
arithmetic       & $\mathbf{51.5\%}$ & $47.6\%$ \\
all tasks        & $\mathbf{60.7\%}$ & $58.8\%$ \\
\bottomrule
\end{tabular}
\end{center}

The reversal is sharpest exactly where the isolated gain was largest: \textbf{mul} goes from unfrozen's
$+9$-point isolated win to a 16-point \emph{loss} under shared load (frozen $59\%$ vs unfrozen $43\%$).
Letting the identity drift to specialise for one operation corrupts the residual every other task
depends on; in isolation there is nothing to corrupt, on a mixed corpus there is. So freezing is not a
per-task generalisation trick (the codec-seeded \emph{initialisation} is what generalises), but under
the realistic shared-corpus regime it is also \textbf{not} a cost: it wins. What it unambiguously buys
on top of that is legibility: because the identity stays on the codec geometry, the answer decodes from
the frozen bands with no probe (face-decode column above, and Section~\ref{sec:inside}). I report the
frozen configuration as the headline on both grounds. Reproducible via
\code{dotnet run --project bench -c Release -- --columnar --tuned-baseline} (isolated) and
\code{dotnet run --project bench -c Release} (shared-corpus ablation, printed under the main table).

\subsection{Reading the computation from the inside}\label{sec:inside}

Because every activation is a phasor face, PrismFormer's internal state is directly \emph{decodable}; a
dense transformer's is not. I test this on the two operations with a codec homomorphism, single-digit
addition (linear band) and multiplication (log band), training both models to convergence and then
inspecting. On raw single-digit accuracy the two models are comparable (single-digit tables are small
and memorisable; the transformer is competitive or better), so the point here is not accuracy but
\emph{legibility}. Decoding the final-position face against the codebook, with \textbf{no trained
probe}, recovers the answer from the frozen identity bands as well as the model's own readout does:

\begin{center}
\begin{tabular}{@{}lcc@{}}
\toprule
decode of the final face (no probe) & add & mul \\
\midrule
homomorphic band only              & 35.2\% & 34.4\% \\
full frozen identity (both bands)  & 45.5\% & 75.0\% \\
model's own readout                & 45.5\% & 75.0\% \\
\bottomrule
\end{tabular}
\end{center}

Identity $=$ readout: the answer literally lives in the codec geometry of the activation. The only way
to read the dense transformer is to \emph{fit} a linear probe to its hidden state
($75.0\% \pm 22.7\%$ add, $62.5\% \pm 13.4\%$ mul), and even that is a trained decoder, not a property
of the representation. That contrast, free and probe-free legibility versus a fitted probe, is the
concrete pay-off of phasor faces. Run \code{dotnet run --project bench -c Release -- --inspect}.

\subsection{Language modelling}

Character-level next-token modelling on a small embedded public-domain text, at matched size, both
temperature-calibrated on a validation slice:

\begin{center}
\begin{tabular}{@{}lcc@{}}
\toprule
model & next-char accuracy & bits/char (lower better) \\
\midrule
\textbf{PrismFormer} & $\mathbf{41.0\% \pm 1.9\%}$ & $\mathbf{3.399 \pm 0.081}$ \\
transformer          & $21.7\% \pm 3.6\%$          & $4.652 \pm 0.273$ \\
unigram baseline     & $19.1\%$                    & $5.322$ \\
\bottomrule
\end{tabular}
\end{center}

Both models beat the unigram baseline (the transformer on bits/char, so it is learning, not broken),
and PrismFormer's char-face embeddings win decisively at this size: the algebraic architecture is
competitive at language, not only arithmetic. Both are weak in absolute terms at $\approx 10^5$
parameters on 1{,}670 characters; the claim is the matched-size comparison, not an absolute language
result. Run \code{dotnet run --project bench -c Release -- --lm}.

\subsection{Exact-merge verification}\label{sec:exactmerge}

Splitting a batch into $K$ shards, computing them sequentially and in parallel, and reducing in a fixed
index order yields gradient tensors with an identical bitwise (FNV) checksum
(\code{0x945A7A4E8E0B9F3B}), and a re-run reproduces it exactly. The \code{verify/} kit reproduces this
in seconds, at the production dimensions. This is the property that lets a mismatched-CPU colony sum into
one coherent model.

A checksum verifies the reduction; a small demonstration shows what it buys, across genuinely separate
model instances. I instantiate a master and four \emph{separate} models and give each slave the data for
only one of four toy skills (copy the first token, echo the last, max, min). Each round every slave
computes the gradient for its one skill, serialises it, the master sums the four shards (\code{Grads.Add})
and takes a single step, and the merged step is copied back to all.

\begin{center}
\begin{tabular}{@{}lcccc@{}}
\toprule
model & copy & last & max & min \\
\midrule
master (after merged rounds)                     & $100\%$ & $100\%$ & $99\%$ & $99\%$ \\
a slave that only sent \emph{copy} gradients     & $100\%$ & $100\%$ & $99\%$ & $99\%$ \\
control trained on \emph{copy} alone             & $100\%$ & $3\%$   & $12\%$ & $26\%$ \\
\bottomrule
\end{tabular}
\end{center}

The slave that only ever computed the \emph{copy} gradient ends up fluent in all four skills and
byte-identical to the master, while a model trained on \emph{copy} alone stays near chance on the other
three: it acquired skills it never saw data for, purely from the merged step. And the master's weights are
\textbf{bit-identical} (max absolute parameter difference $0.0$) to a single-process model trained on the
union of the four shards in the same order, so the four-way split is lossless, not approximate. This is
the merge property made concrete across separate instances, not just a checksum.

The merge holds along the whole training trajectory, not only at the end. Running the same shards as a
merged swarm and as a single-process reference in lockstep, the swarm's loss, held-out accuracy, \emph{and
its parameters} match the reference at every checkpoint, to the bit:

\begin{center}
\begin{tabular}{@{}rccc@{}}
\toprule
round & swarm loss / single-proc & swarm acc / single-proc & max param diff \\
\midrule
0   & $3.428$ / $3.428$ & $38.1\%$ / $38.1\%$ & $0.0$ \\
100 & $0.870$ / $0.870$ & $74.8\%$ / $74.8\%$ & $0.0$ \\
300 & $0.018$ / $0.018$ & $98.6\%$ / $98.6\%$ & $0.0$ \\
600 & $0.007$ / $0.007$ & $99.8\%$ / $99.8\%$ & $0.0$ \\
\bottomrule
\end{tabular}
\end{center}

Loss falls to a floor and accuracy rises to $99.8\%$ over $600$ merge rounds, and the two columns are the
same number in every row: the swarm run is not merely close to single-machine training, it \textbf{is}
single-machine training step for step, because a lossless gradient sum has no drift to accumulate. And it
survives leaving the process: with one host and several worker cells over real (loopback) TCP sockets, the
swarm's loss tracks a single-node reference to four decimals at every round ($3.46\to2.08$), the final
model is bit-identical to that reference, and every worker cell is byte-identical to every other. So the
merge is demonstrated in-process, across separate instances, along the full convergence curve, and over
sockets. What is left is the genuinely \emph{distributed} systems claim, a geo-distributed run across
mismatched unreliable machines with a wall-clock scaling curve, and a measured comparison against ordinary
asynchronous approximate merging, which stay future work (Section~\ref{sec:lim}). Run
\code{dotnet run --project bench -c Release -- --swarmdemo} (and \code{--swarmconverge}), and
\code{PrismNet swarmtest} for the over-sockets version.

\subsection{The codec-seeded baseline control}\label{sec:codecbaseline}

Because PrismFormer's number embeddings are seeded from the phasor codec while the transformer baseline
starts random, a fair reader asks whether the head-to-head gap is \emph{initialisation} rather than
architecture. I test the confound at its source: hand the transformer the \emph{same} codec seeding and
re-measure. Three models train on the same in-range held-out add/sub split (operands 0--12, 20\% of pairs
held), over $5$ seeds: AlgFormer with codec-seeded embeddings; a transformer with random embeddings; and
\emph{the same transformer with its embeddings seeded from the identical codec}, so the only difference
from the random one is the initialisation. The transformer is given a warmup schedule so that it fits
train, and at $\sim$805k parameters it is $3.3\times$ larger than AlgFormer's $\sim$242k.

\begin{center}
\begin{tabular}{@{}lccc@{}}
\toprule
model & params & train & held-out (mean$\pm$sd, 5 seeds) \\
\midrule
AlgFormer (codec-seeded)   & 242k & $100\%$ & $45.8\% \pm 15.9\%$ \\
transformer (random init)  & 805k & $95\%$  & $31.3\% \pm 21.3\%$ \\
transformer (codec-seeded) & 805k & $63\%$  & $\mathbf{9.1\% \pm 6.9\%}$ \\
\bottomrule
\end{tabular}
\end{center}

The confound is refuted, and in the strongest way: seeding the transformer from the codec does
\textbf{not} help it, it \textbf{hurts} it, on every one of the five seeds. Held-out falls from $31.3\%$
to $9.1\%$, and the seeded transformer cannot even fit the training data as well ($63\%$ vs $95\%$). So
the codec is not a transferable free head-start that inflates AlgFormer's numbers; the same phasor
geometry that AlgFormer's bind and correlate operations are built to read is a \emph{liability} to
attention, which cannot exploit it. I am deliberately not over-reading the rest of the table: on this
small single-token task AlgFormer and a tuned random-init transformer are within seed-to-seed noise
($45.8\pm15.9$ vs $31.3\pm21.3$, and the transformer wins on two of the five seeds), so the clean
architectural separation is the one shown on the worked multi-digit and relational tasks
(Sections~\ref{sec:worked} and~\ref{sec:relational}), not here. What this control establishes is narrower
and cleaner: the initialisation confound does not explain the gap, because the initialisation does not
transfer. Run \code{dotnet run --project bench -c Release -- --codec-baseline}.

\subsection{Skill-sharing by continuous coupling}\label{sec:skillshare}

The swarm's other half is not just that gradients \emph{merge} but that skill \emph{spreads}: a node
should pick up what its peers have learned. I test this in a single process, simulating the mesh's weight
coupling directly (a live multi-node run is still future work, Section~\ref{sec:lim}). The setup is the
elastic-averaging move the mesh actually uses, a continuous pull of each node's weights toward a running
average of its peers' \parencite{zhang2015easgd,mcmahan2017fedavg}, on a small multi-skill arithmetic
corpus where different nodes are trained on different skills. Three findings, and I am careful about the
boundary of each.

First, \textbf{one-shot averaging of independently-trained models fails.} Train two models to competence
and average their weights once: with a shared initialisation and the frozen codec it holds ($72\%$), but
with \emph{different} initialisations it collapses to near-chance ($9\%$ with the codec, $6\%$ without).
This is the linear-mode-connectivity / permutation-symmetry barrier \parencite{ainsworth2023rebasin}: two
independently-initialised networks sit in different basins, and their midpoint is not a good model. The
codec alone does not rescue it, because the codec pins the \emph{embedding} coordinate system, not the
relation-bank weights.

Second, \textbf{continuous coupling bridges different initialisations where one-shot averaging cannot.}
Replace the single average with a gentle pull-to-mean applied \emph{throughout} training (elastic
averaging): two different-init nodes now reach $\mathbf{75\%}$ without the codec and $\mathbf{99\%}$ with
the frozen codec, while their weights stay measurably distinct (divergence $0.07$ and $0.02$). Skill
transfers directly: a node picks up a peer's competence purely through weight coupling, with no shared
training data.

Third, the \textbf{codec is a bounded enhancer, not a magic key.} Coupling already works on a plain
network ($75\%$); the frozen codec lifts it to near-lossless ($99\%$) by giving every node the same fixed
coordinate system, so their updates stay commensurable instead of drifting into incompatible bases. The
honest claim is therefore precise: continuous elastic coupling shares skill across
independently-initialised nodes on \emph{any} network, and the algebraic codec makes that sharing
near-lossless. It is neither ``only the codec can do this'' (a plain net reaches $75\%$) nor a new
optimiser (elastic averaging is standard \parencite{zhang2015easgd}); it is a measured, specific benefit
of a shared representational coordinate system. Averaging weights from a \emph{shared} initialisation is
already known to work \parencite{wortsman2022soups}; the contribution here is that the fixed codec extends
near-lossless sharing to nodes that never shared one. Run
\code{dotnet run --project bench -c Release -- --average}. The companion \code{--swarmdemo} shows the same
bleed as a before/after matrix: two shared-init specialists (copy at $100\%$ / max at $12\%$, and the
mirror) elastic-average into two generalists that each gain the other's skill ($12\to94\%$ and
$25\to88\%$) as their weights collapse together (divergence $0.27\to0.008$).

That consensus has a trajectory, and it is worth laying beside the exact-merge curve, because the two
swarm mechanisms converge in \emph{different} senses. Coupling the two specialists and logging over rounds,
the divergence between them falls to a small floor while each node's \emph{other}-skill accuracy climbs:

\begin{center}
\begin{tabular}{@{}rccc@{}}
\toprule
round & divergence(A,B) & A's other-skill & B's other-skill \\
\midrule
0   & $0.270$ & $12.5\%$ & $25.0\%$ \\
300 & $0.031$ & $27.1\%$ & $59.0\%$ \\
600 & $0.009$ & $84.6\%$ & $44.9\%$ \\
900 & $0.008$ & $93.8\%$ & $87.9\%$ \\
\bottomrule
\end{tabular}
\end{center}

The contrast with the gradient-sum curve (Section~\ref{sec:exactmerge}) is the point. Gradient summing is a
lossless \emph{sum}, so the merged run is bit-identical to single-machine (max param diff $0.0$); weight
averaging is a lossy \emph{consensus}, so the divergence shrinks toward a small \emph{nonzero} floor
($0.008$, not $0$) while skill bleeds across ($12.5\to93.8\%$ and $25\to87.9\%$). Both converge: one to the
same bits, the other to a shared, near-lossless middle. Run
\code{dotnet run --project bench -c Release -- --avgconverge}.

\subsection{Reverse inference, recovered}\label{sec:revinfer}

The reverse-inference failure of Section~\ref{sec:relational} (train ``opposite of A = B'', hold out
``opposite of B = A''; $0\%$ for both models) deserves a closer look, because the codec's own machinery
addresses it. The right diagnosis is that the $0\%$ is a \emph{directional-recall} failure, not a
reasoning failure: a closed-book model trained on $A\!\to\!B$ stores a directed map, and nothing in it
says the relation runs backward. Reversing, once you have the fact, is trivial. I show this three ways on
a set of symmetric pairs (run \code{dotnet run --project bench -c Release -- --revinfer}):

\begin{center}
\begin{tabular}{@{}lcc@{}}
\toprule
 & forward & reverse (held-out) \\
\midrule
closed-book, trained (the setup above)         & $100\%$    & $0\%$ \\
open-book scratchpad (fact in context)         & ---        & $\mathbf{100\%}$ \\
algebra (commutative bind, no training), 12 / 24 pairs & $100/83\%$ & $\mathbf{100/83\%}$ \\
\bottomrule
\end{tabular}
\end{center}

\textbf{Scratchpad.} Provide the forward fact in context and the reverse becomes ``select the element
that is not the query'', the same copy operation the model already does at $100\%$
(Section~\ref{sec:relational}); it generalises to unseen pairings of seen tokens at $100\%$. This does
not conjure a fact taught in only one direction; it removes the \emph{recall} bottleneck by making the
fact available, which is exactly what a retrieval step or a reasoning scratchpad provides.

\textbf{Algebra.} Store a symmetric relation as a \emph{commutative} bind $L \bind R$, retrieved by
unbinding the query ($L\bind R\bind R^{-1} = L$, and the bind commutes), so forward and reverse are the
\emph{same} operation and the reverse is free, with no training. The only shortfall from $100\%$ is HRR
crosstalk as the associative bundle fills, a capacity property rather than an asymmetry. This is the
arithmetic-as-algebra move of Section~\ref{sec:codec} applied to relations: the reverse falls out of the
codec the way a sum or product does.

So reverse inference is not a wall. It is $0\%$ only in the closed-book, directional-recall regime, and
either a scratchpad (retrieval) or an algebraic encoding recovers it. I still report the closed-book
$0\%$ as the honest number for the end-to-end trained model, and treat this as the constructive next step.

\section{Limitations}\label{sec:lim}

I state these plainly. (i) \textbf{Arithmetic scope.} I claim generalisation to unseen operand
\emph{combinations within} the trained range (Section~\ref{sec:worked}, worked multi-digit add/sub and
multiply/divide by a digit). I do \textbf{not} claim: \emph{out-of-range magnitude} extrapolation
(operands larger than any seen is $\approx$1\% for both models, because the end-to-end model learns the
mapping statistically, not via the codec's exact homomorphism); \emph{length} extrapolation (more
digits than trained, a hard positional-generalisation problem I do not target); or full multi-digit
$\times$ multi-digit multiplication, which needs a partial-product scratchpad
(Section~\ref{sec:worked}). (ii) \textbf{Reverse-inference} (inferring a held-out symmetric/reverse
relation from the trained direction) is \textbf{0\% for both models} at this scale
(Section~\ref{sec:relational}) in the closed-book, end-to-end setting; it is recoverable with a scratchpad
(retrieval) or an algebraic commutative-bind encoding (Section~\ref{sec:revinfer}), which I treat as the
constructive next step rather than a solved integration. (iii) \textbf{Freezing is regime-dependent}: on a single operation in
isolation an unfrozen (still codec-seeded) identity generalises better; freezing only wins under shared
multi-task load, and its unconditional benefit is legibility (Section~\ref{sec:worked}). I keep it as
the headline for the shared-corpus and legibility reasons and state the isolated-task caveat. (iv)
\textbf{Scale.} Results are at $\approx 10^5$ parameters; large-scale convergence of the full colony is
demonstrated to \emph{run} but not yet to \emph{converge} at scale. (v) \textbf{Novelty.} The phasor codec and RFF
encoding are prior art, and neither ``decodable numbers'' (FoNE) nor ``algebra in the embedding'' (NIF)
is first (Section~\ref{sec:related}); my contribution is the specific combination, one fixed
complex-multiply bind over linear and log bands, decodable by correlation and unified across words and
numbers, plus the relation-banks and the exactly-mergeable substrate and the swarm it enables.

\section{Future work and open questions}

I would rather name the sharpest objections to this work than let a reader find them unaided. Each is a
concrete next step, not a rhetorical hedge.

\textbf{The codec-seeded baseline confound, now controlled.} PrismFormer's numeric embeddings are seeded
from the phasor codec while the transformer baseline starts random, which could confound architecture
with initialisation. I ran the control (Section~\ref{sec:codecbaseline}): handing the transformer the
same codec seeding does not help it and in fact hurts it (held-out $31.3\%\to9.1\%$ over five seeds, and
it fits train less well), so the codec is not a transferable head-start that explains the gap. What that
control does \emph{not} settle, because this single-token task is small and high-variance, is a clean
architecture-versus-transformer margin on it; that separation lives in the worked multi-digit and
relational results, and a stronger baseline is discussed next.

\textbf{A tuned baseline does not close the gap.} My main transformer baseline is a hand-derived,
gradient-checked encoder without layer normalisation, matched by parameter count. It fits its training
data, so it is not a strawman, but it is not a modern reference either, so I re-ran the fair multi-task
suite against a properly tuned transformer: pre-norm layer normalisation (its backward hand-derived and
gradient-checked to $0$), a warmup-then-cosine schedule, and a tuned Adam, still parameter-matched
($298{,}752$ vs $295{,}936$, five seeds). It reaches $50.8\%$ held-out on the fair tasks, against
$49.1\%$ for the untuned baseline and $65.9\%$ for PrismFormer: tuning the baseline moves it under two
points and does not close the gap. What I still do not compare against is existing complex-valued and
VSA / neuro-symbolic models, which is needed before the architectural claim is fully settled.

\textbf{The swarm: property and convergence curve shown, distribution at scale not.} I show that sharded
gradients merge bit-for-bit, that the merged run converges bit-identically to single-machine training
along the whole curve, in-process and over loopback TCP (Section~\ref{sec:exactmerge}), and that skill
spreads across independently-initialised nodes through continuous coupling (Section~\ref{sec:skillshare}).
What every experiment shares is that it runs on one machine. A genuine \emph{geo-distributed} run across
mismatched, unreliable machines, with a wall-clock scaling curve and a fault-injection test at real
network scale, is future work; until it exists, the swarm should be read as a demonstrated
\emph{merge-and-share property with a matching convergence curve} plus an \emph{architecture}, not a
demonstrated at-scale distributed-training result. The value of bit-exactness over ordinary asynchronous
approximate merging is likewise something I argue for but have not yet measured.

\textbf{The homomorphism does not extrapolate, and that narrows the thesis.} Within the trained range
the codec makes arithmetic decodable and the model generalises to unseen operand combinations; out of
range it collapses toward chance, and (Section~\ref{sec:lim}) it is the network learning a statistical
mapping on the codec's representation, not the codec's exact homomorphism, that does the generalising.
So the honest thesis is narrower than ``arithmetic as algebra'': the codec supplies a decodable,
well-conditioned representation and initialisation, and the network learns on top of it. Making the
exact homomorphism itself carry out-of-range values and multi-digit-by-multi-digit multiplication (via a
partial-product scratchpad, Section~\ref{sec:worked}) is open.

\textbf{Scale and statistics.} Everything here is at roughly $10^5$ parameters with 4 to 8 seeds, and
several per-task standard deviations are large relative to their means; I report means and standard
deviations but not confidence intervals or significance tests across the task suite. Whether any of this
survives one or two orders of magnitude more parameters is unknown, and the absence of layer
normalisation, which currently caps trainable depth, is the first thing that has to change to get there.

None of this is settled, and I would rather ship it saying so.

\section{Conclusion}

I set out to build a representation that reasons instead of recalls, and the honest result is narrower
and more concrete than that ambition: numbers written as phasors make arithmetic something you can
decode straight out of the representation, and a model small enough to live on one CPU, with gradients
that merge to the bit, makes a fault-tolerant training swarm possible. Those two things came from the
same idea and I think they stand. I have tried to make every claim here something you can run and check
for yourself, including one of my own that the evidence knocked down, and the verification kit and
per-section commands are there so anyone can reproduce, or refute, each number.

\section*{Reproducibility}

The exact code that produced every number above is bundled with this archive as a frozen source
snapshot (repository tag \texttt{paper-v3}), so it does not depend on the live repository, which
continues to evolve. Build with the .NET 8 SDK; each command below is run from the snapshot root.

\begin{itemize}
\item \code{verify/}: the arithmetic-from-the-codec and exact-merge checks in seconds, no training
(\code{dotnet run --project verify/Verify -c Release}).
\item \code{bench/}: the parameter-matched head-to-heads, the default run (Section~\ref{sec:relational}),
plus \code{--columnar} (Section~\ref{sec:worked}, with the frozen-identity ablation), \code{--inspect}
(Section~\ref{sec:inside}, mechanistic), \code{--lm} (\S\,4.5), \code{--codec-baseline}
(Section~\ref{sec:codecbaseline}), \code{--average} (Section~\ref{sec:skillshare}), and \code{--revinfer}
(Section~\ref{sec:revinfer}).
\item \code{PRIOR\_ART.md}: a timestamped defensive-publication disclosure of the methods.
\end{itemize}

\section*{Tool and computational resource disclosure}

In the spirit of the Leiden Declaration on Artificial Intelligence and Mathematics, I disclose the tools
behind this work. The ideas, the design, and the direction are mine: the axioms and Platonic-space
framing, the phasor codec and its homomorphism, the relation-bank transformer, and the exactly-mergeable
swarm, together with every design decision and the choice of what to test. I used a large-language-model
coding assistant as an \emph{execution} tool: to implement the software and benchmark suite, to run and
tabulate experiments, and to draft and edit this manuscript to my specification and under my review. It
carried out my instructions; it did not originate the research. Every claim, number, and argument was
checked and verified by me, and I retain full responsibility for the correctness and adequacy of the
results; where the evidence contradicted one of my own hypotheses, I report the contradiction
(Section~\ref{sec:worked}). No claim here rests on an unverified automated argument: every quantitative
result is produced by the reproducibility kit (frozen snapshot, tag \texttt{paper-v3}) and can be re-run
or refuted independently. All experiments ran on commodity CPUs; no result depends on specialised
accelerators.

\printbibliography[title={References}]

\vspace{1em}
\emph{Please open an issue if any attribution is incomplete or mis-stated; correct credit to prior work
matters to me.}

\end{document}
